\chapter{Andre Beteille on Caste}

\section{The Segmentary System}

\subsection{The Two Books and the Weberian Frame}

\begin{KeyConcept}
\begin{itemize}
\item Andre Beteille wrote two landmark books: \textbf{'Caste, Class and Power: Changing Patterns of Stratification in a Tanjore Village'} and \textbf{'Caste, Old and New'}.
\item The title 'Caste, Class and Power' echoes \textbf{Max Weber} (status, class, power) --- Beteille replaces \textbf{'status' with 'caste'} (caste represents one's status in the hierarchy).
\item Beteille helped European scholars and Indian readers alike understand caste.
\end{itemize}
\end{KeyConcept}

\subsection{Shripuram: The Three Clusters}

\begin{ExampleBox}
\begin{itemize}
\item Beteille did field work in \textbf{Shripuram} village, \textbf{Tanjore} district, Tamil Nadu --- on the advice of \textbf{M. N. Srinivas}, who told him \emph{not} to study West Bengal (Beteille is a Bengali Brahmin; his own biases would prevent objective data) but to go somewhere completely new.
\item The village has \textbf{three clusters}: \textbf{Agraharam} (Brahmins), \textbf{Kudiyana streets} (non-Brahmins), and \textbf{Cheri} (Adi Dravidas / untouchables).
\item Population order: \textbf{non-Brahmins > Brahmins > Adi Dravidas} (non-Brahmins are the majority).
\item There is \textbf{complete segregation} based on \textbf{purity and pollution}: a Brahmin cannot pass through the untouchables' streets (he would need a ritual bath); the untouchables too regard a Brahmin's presence in Cheri as profane. This is \textbf{social distancing based on caste}.
\end{itemize}
\end{ExampleBox}

\subsection{The Five Brahmin Categories (and their sub-divisions)}

\begin{KeyConcept}
\begin{itemize}
\item Beteille divided the Brahmins of Shripuram into \textbf{five categories}: (1) \textbf{Kannada Madhva}, (2) \textbf{Telugu Smartha}, (3) \textbf{Tamil Smartha}, (4) \textbf{Tamil Sri Vaishnava}, (5) \textbf{Temple Priests}.
\item \textbf{Smartha} Brahmins are followers of \textbf{Shankaracharya (Adi Shankar)} --- \textbf{Advaita (monism)} --- and worship both Shiva and Vishnu. \textbf{Sri Vaishnava} Brahmins are followers of \textbf{Ramanuja} --- \textbf{Vishishtadvaita (qualified monism)} --- and worship Vishnu (Venkateshwara) only.
\item These are \textbf{sectarian/doctrinal differences within the same caste} (Brahmins), not merely differences of whom they worship.
\item \textbf{Temple Priests} are a different league of Brahmins (not every Brahmin is a temple priest).
\end{itemize}
\end{KeyConcept}

\begin{KeyConcept}
\begin{itemize}
\item The categories are further subdivided: the \textbf{Telugu Smartha} into Kona Seema, Velnadu, Mulhanadu; the \textbf{Tamil Smartha} into Vadama, Brihacharanam, Ashtasahasram and Vattiman (and Vadama into Vadadesha and Chozhadesha, Brihacharanam into Kandramanikya, etc.); the \textbf{Tamil Sri Vaishnava} into Vadagalai and Thengalai; and the \textbf{Temple Priests} into Bhattachar and Kurukkal.
\item The point: \textbf{Brahmins are not homogeneous} --- there are segments within segments (fission within Brahmins).
\end{itemize}
\end{KeyConcept}

\subsection{Non-Brahmins and Adi Dravidas}

\begin{KeyConcept}
\begin{itemize}
\item The \textbf{non-Brahmins} are the people of the \textbf{DMK / Self-Respect / anti-Brahmin movement} (led by Periyar, who was the guardian of the untouchables). They are divided into three kinds of castes:
\item \textbf{Artisan castes} (potter, goldsmith), \textbf{Servicing castes} (barber, musicians), and \textbf{Cultivating castes} (peasants): \textbf{Vellalans}, \textbf{Padayachi}, \textbf{Gaundans} --- the \textbf{Vellalans} are the majority and are traditional cultivators working on Brahmin land (the Brahmins are \textbf{landlords / mirasdars}).
\item The \textbf{Adi Dravidas} (untouchables) are divided into \textbf{Pallans} (Hindu untouchables) and \textbf{Parayans} (Christian untouchables). 'Parai' is a musical instrument played at death processions --- associated with the impure.
\end{itemize}
\end{KeyConcept}

\subsection{Caste as a Segmentary System}

\begin{KeyConcept}
\begin{itemize}
\item \textbf{Caste is a segmentary system} --- consisting of multiple segments (Brahmin, Smartha, Vadama, Vadadesha... are all 'jatis').
\item \textbf{The principle of segmentation is constant, but the criteria for segmentation may vary.}
\item Before Beteille, people confused caste and sub-caste; Beteille showed these are all \textbf{jatis activated in specific contexts} --- one thing, defined by context.
\end{itemize}
\end{KeyConcept}

\subsection{Elastic Ethnicity and Context}

\begin{KeyConcept}
\begin{itemize}
\item Beteille uses \textbf{Fredrik Barth's} idea of \textbf{elastic ethnicity / elastic boundaries}: which identity (jati) you activate depends on context.
\item Example: 'Where are you from?' --- New Rajendra Nagar, Central Delhi, Punjabi, or Indian? All are identities (all are 'jatis'); which one you use depends on context (against Pakistan in cricket you are Indian; in the IPL you are Delhi).
\item The same applies to the divisions of caste: in the \textbf{context of marriage}, the local sub-sub-caste identity becomes important; in \textbf{sectarian/doctrinal} contexts, the Smartha/Sri Vaishnava identity; in \textbf{political} contexts, the Brahmin identity.
\end{itemize}
\end{KeyConcept}

\subsection{The Four Levels of Segmentation}

\begin{KeyConcept}
\begin{itemize}
\item Beteille's \textbf{four levels of segmentation} in Shripuram:
\item \textbf{Level 1} (clusters: Brahmin/non-Brahmin/Adi Dravida) --- activated in the \textbf{political context}.
\item \textbf{Level 2} (sectarian: Smartha vs Sri Vaishnava) --- activated in the \textbf{sectarian/doctrinal context}.
\item \textbf{Levels 3 and 4} (sub-sub-castes: Vadadesha vs Chozhadesha, etc.) --- activated in the \textbf{context of marriage}.
\item In the marriage context a Brahmin is not just a Brahmin --- he is Smartha or Sri Vaishnava, and within Smartha, Vadadesha cannot marry Chozhadesha. In the sectarian context he is Smartha vs Sri Vaishnava. In the political context, Smartha and Sri Vaishnava \textbf{unite as 'Brahmin'} (against non-Brahmin movements) --- 'a man is Smartha only in relation to Sri Vaishnava'.
\end{itemize}
\end{KeyConcept}

\begin{ExampleBox}
\begin{itemize}
\item In the \textbf{1962 election} in Shripuram, the candidate \textbf{Siyar Patta Bhimaran} was a Smartha Brahmin fighting the non-Brahmins; Brahmins (including Sri Vaishnavas) voted for him --- his Smartha identity became secondary, and the \textbf{Brahmin identity became primary} in the political context.
\item This illustrates Srinivas's principle (Beteille is Srinivas's student): \textbf{context defines status}.
\end{itemize}
\end{ExampleBox}

\subsection{Caste, Class and Power: Cumulative Inequality in Traditional India}

\begin{KeyConcept}
\begin{itemize}
\item In \textbf{traditional Shripuram} (before colonial policies), there was a \textbf{vertical hierarchy}: Brahmins at the top (ritual hierarchy) and Adi Dravidas at the bottom --- but there was also a \textbf{secular hierarchy based on wealth and power}.
\item Traditional India: \textbf{Brahmins were landlords AND political elites}; non-Brahmins were \textbf{cultivators on Brahmin land AND politically subordinate}; Adi Dravidas were \textbf{landless labourers}.
\item Therefore \textbf{caste was class} --- from your caste came your class, your occupation, your wealth and your power.
\item This is \textbf{cumulative inequality} (Beteille's term): caste, status, wealth and power are \textbf{tightly integrated} --- the Brahmins were rank 1 in status, wealth and power (1-1-1); the non-Brahmins rank 2, the Adi Dravidas rank 3.
\item You could not change your occupation because you could not change your caste.
\end{itemize}
\end{KeyConcept}

\begin{KeyConcept}
With the \textbf{advent of Westernization/colonial policies}, the \textbf{forms of stratification began to change} in Tanjore village (the 'changing patterns' of the book's subtitle) --- the detailed discussion follows.
\end{KeyConcept}

\begin{CriticalPoint}
A related note: in Indian sociology there is \textbf{no hard distinction between sociology and anthropology} --- \textbf{Satish Deshpande, Patricia Uberoi and Nandini Sundar} call it a \textbf{'metropolitan division of labour'} (sociologists study urban Western societies; anthropologists study rural/tribal societies) --- which is nothing but the \textbf{politics of knowledge}. Beteille used field work, which sociologists do too.
\end{CriticalPoint}

\begin{QuickRevision}
\begin{itemize}
\item Two books: 'Caste, Class and Power' (Tanjore/Shripuram) and 'Caste, Old and New'; title echoes Max Weber (status$\rightarrow$caste).
\item Shripuram: Agraharam (Brahmins), Kudiyana (non-Brahmins), Cheri (Adi Dravidas); non-Brahmins > Brahmins > Adi Dravidas; complete segregation (purity/pollution).
\item Five Brahmin categories: Kannada Madhva, Telugu Smartha, Tamil Smartha, Tamil Sri Vaishnava, Temple Priests; Smartha (Shankaracharya, Advaita) vs Sri Vaishnava (Ramanuja, Vishishtadvaita).
\item Non-Brahmins: artisan/servicing/cultivating castes (Vellalans majority); Adi Dravidas: Pallans (Hindu) vs Parayans (Christian).
\item Caste = segmentary system; principle of segmentation constant, criteria vary; Brahmins not homogeneous (fission).
\item Elastic ethnicity (Barth): which jati activates depends on context (marriage $\rightarrow$ sub-sub-caste; sectarian $\rightarrow$ Smartha/Sri Vaishnava; political $\rightarrow$ Brahmin).
\item Four levels: L1 political (Brahmin), L2 sectarian (Smartha/Sri Vaishnava), L3-4 marriage (Vadadesha/Chozhadesha); 'context defines status'.
\item Traditional India: caste = class (Brahmins = landlords + political elites); cumulative inequality (Beteille); westernization begins changing stratification.
\end{itemize}
\end{QuickRevision}

\section{Caste, Class and Power}

\subsection{Segments as Status Groups (Weber)}

\begin{KeyConcept}
\begin{itemize}
\item Caste is a \textbf{segmentary system} whose segments (Brahmin, Smartha, Vadama, etc.) are \textbf{status groups} in themselves --- in Weber's sense.
\item A \textbf{status group} is a community marked by a \textbf{distinct lifestyle} and \textbf{endogamy}. \textbf{Max Weber} gave the \textbf{Brahmin} as his example of a status group in India.
\item \textbf{Beteille = Max Weber + field work}: he accepts that Brahmin is a status group, but shows it is \textbf{not homogeneous} --- each segment within it is itself a status group (a jati).
\item So the Smartha will marry only a Smartha (endogamy), not a Sri Vaishnava --- but the Smartha \textbf{can vote for} a Sri Vaishnava, because in the political context they are both Brahmins.
\end{itemize}
\end{KeyConcept}

\subsection{Traditional Shripuram: The Rank Matrix (Cumulative Inequality)}

\begin{KeyConcept}
\begin{itemize}
\item In \textbf{traditional Shripuram}, there were two hierarchies: the \textbf{ritual hierarchy} (Brahmin > non-Brahmin > Adi Dravida) and the \textbf{secular hierarchy} (based on wealth and power).
\item \textbf{Brahmins} were \textbf{landlords (mirasdars)} and \textbf{political elites} (the headman); \textbf{non-Brahmins} were \textbf{cultivators} on Brahmin land and \textbf{political subordinates}; \textbf{Adi Dravidas} were \textbf{landless labourers} (the lowest, like bonded labour).
\item Appointment was based on \textbf{caste, not merit} --- a \textbf{patrimonial} bureaucracy (Weber).
\item This is Beteille's \textbf{rank matrix}: caste, class and power run together (1-1-1, 2-2-2, 3-3-3) --- your power and wealth were \textbf{derived from your status} (Weber), i.e. \textbf{caste = class = power}, tightly integrated (cumulative inequality).
\item (Contrast with Srinivas's Rampura, where the Vokkaliga peasants were dominant; in Shripuram the Brahmins are landlords and elites.)
\end{itemize}
\end{KeyConcept}

\subsection{Economic Modernization: Land, Occupations, Education}

\begin{KeyConcept}
\begin{itemize}
\item \textbf{Economic modernization} includes private property, cash economy, industrialization, new means of production, real estate, and the \textbf{free market economy}.
\item Its effects on Shripuram:
\item 1. \textbf{Land became a marketable commodity} --- with a cash economy land can be bought and sold, so land was \textbf{freed from the clutches of the caste structure} (the Brahmins began loosening their hold over land ownership).
\item 2. \textbf{Rise of new occupations} --- no longer only landlord/cultivator/labourer; now IAS/IPS officers, civil engineers, computer professionals.
\item 3. \textbf{Rise of English-medium education} --- education as human capital (improving skilled, productive labour and hence GDP).
\end{itemize}
\end{KeyConcept}

\subsection{Political Modernization: Power Without Land}

\begin{KeyConcept}
\begin{itemize}
\item \textbf{Political modernization} introduced a \textbf{democratic political structure} (universal adult franchise, political parties, Panchayati Raj, vote banks; the caste panchayat weakened).
\item Its uniqueness: it gives you \textbf{power even without land ownership} --- the basis of power is no longer only land.
\item \textbf{Party politics} replaced village \textbf{caste politics} --- and with parties came \textbf{networking outside the village} (before, a Vadama Brahmin had no contacts outside Shripuram; a BJP member does).
\item This also gave rise to \textbf{elites} --- 'lions and foxes' (manipulators) --- produced by the democratic structure.
\item Impact: (1) \textbf{land ownership became accessible to non-Brahmin cultivators}; (2) \textbf{political power moved to non-Brahmin elites} --- democracy is majority rule, and the non-Brahmins have the numbers.
\item Yet party politics now \textbf{runs caste politics} (parties made on caste lines) --- tradition and modernity coexist (Yogendra Singh; Parsons' pattern variables A and B).
\end{itemize}
\end{KeyConcept}

\subsection{The Changing Rank Matrix: Autonomy of Class and Power}

\begin{KeyConcept}
\begin{itemize}
\item In the traditional rank matrix, the \textbf{degree of correlation between caste, class and power was very strong} (one-to-one). After modernization, the \textbf{correlation is weak} --- the \textbf{degree of variation is strong}.
\item There is a \textbf{shift in the rank matrix}: no hard-line segregation, mobility, and fragmented boundaries.
\item A \textbf{Brahmin need no longer be rank 1 in power and wealth} (though still 1 in status); a \textbf{non-Brahmin can be wealthier and more powerful than a Brahmin}; even \textbf{Adi Dravidas} can be powerful (via reservation and education).
\item \textbf{Class and power gained autonomy from caste} --- the village was undifferentiated before; now \textbf{structural differentiation} begins.
\end{itemize}
\end{KeyConcept}

\subsection{Caste Becoming Heterogeneous: A New Status Group}

\begin{KeyConcept}
\begin{itemize}
\item Beteille's conclusion: \textbf{caste in modern India is becoming heterogeneous} --- in terms of \textbf{income, occupation and education}.
\item A \textbf{new status group is emerging based on achieved statuses} (income, education), replacing the old status group based on \textbf{ascribed features (birth)}.
\item This is also a \textbf{status group} --- its members have a \textbf{distinct lifestyle} (e.g. the \textbf{Indian Foreign Service} officers, with their distinct elitism, manners and social-club etiquette).
\item A \textbf{new caste is developing} in the Indian village: one for which \textbf{ritual knowledge is no longer important} --- it has been replaced by \textbf{English and etiquettes}.
\end{itemize}
\end{KeyConcept}

\begin{ExampleBox}
Example: a Brahmin typist who daily sees an Adi Dravida IAS officer giving him orders feels humiliated --- yet may even want his daughter to marry that officer. Change comes through such marriages: the upper-caste man marrying the lower-caste woman (or vice versa) shows caste becoming \textbf{heterogeneous}.
\end{ExampleBox}

\begin{CriticalPoint}
\textbf{Modernization gives meritocracy, free market, English-medium education and egalitarianism --- but it is also creating a new hierarchy.} It is no longer caste-based; it is based on \textbf{who has adopted westernization, and to what extent}.
\end{CriticalPoint}

\subsection{Modernistic vs Modernizing}

\begin{KeyConcept}
\begin{itemize}
\item Beteille distinguishes two kinds of people (drawing on Srinivas's 'externally vs internally westernized'):
\item \textbf{Modernistic} (externally westernized) --- those who have \textbf{appropriated Western symbols/outlook} (dress, food, appearance) to look and feel Western.
\item \textbf{Modernizing} (internally westernized) --- those who \textbf{promote Western values/thinking} (not merely Western appearance).
\item A democracy needs \textbf{internally westernized} people, but India has more and more \textbf{externally westernized} people --- they look Western but 'speak Hindu' --- and this is \textbf{a threat to democracy}.
\item Beteille hints that \textbf{some Indian elites are in danger of becoming more modernistic than modernizing}.
\end{itemize}
\end{KeyConcept}

\subsection{Conclusion: The Paradox of Modernization}

\begin{KeyConcept}
Beteille notes \textbf{the paradox of modernization in India}: it has \textbf{given birth to new castes} --- castes \textbf{imported directly from the West} (e.g. the IFS officer), along with the \textbf{Western ideology of an open society}.
\end{KeyConcept}

\begin{QuickRevision}
\begin{itemize}
\item Caste = segmentary system; segments are status groups (Weber: Brahmin as status group); Beteille = Weber + field work; Smartha marries Smartha (endogamy) but votes for Sri Vaishnava (politics).
\item Traditional Shripuram rank matrix: Brahmin = landlord (mirasdar) + political elite; non-Brahmin = cultivator + subordinate; Adi Dravida = landless labourer; caste = class = power (cumulative inequality); patrimonial appointment.
\item Economic modernization: land marketable (freed from caste), new occupations, English-medium education (human capital).
\item Political modernization: democratic structure $\rightarrow$ power without land; party politics $\rightarrow$ networking; elites (lions/foxes); land to non-Brahmins, power to non-Brahmin elites (numbers).
\item Rank matrix shifts: weak correlation, strong variation; class/power autonomous from caste; structural differentiation.
\item Caste becomes heterogeneous (income/occupation/education); new status group on achieved statuses; ritual $\rightarrow$ English/etiquettes; new caste (IFS).
\item Modernistic (externally westernized) vs modernizing (internally westernized) --- threat to democracy.
\item Paradox: modernization creates new (imported) castes + open society ideology.
\end{itemize}
\end{QuickRevision}

